\begin{theoremEnd}[combined main,category=positive]{corollary}[Positivity in the Laplacian expansion]
\CombinedResultCheck{cor:positive}{corollary}\label{cor:positive}
Under the finite Gaussian model in \cref{sec:finite-model}, suppose
$m\geq1$ and $I\subset(c-L,c+L)$. Set
$h_0=1$, $\boldsymbol g=(a_1,\ldots,a_m)^T$, and
$f_{\boldsymbol g}=\sum_{j=1}^ma_j\phi_{j,L}$.
Then $f_{\boldsymbol a}=f_{\boldsymbol g}+a_0h_0$, with inverse
obtained by reinserting $a_0$.

For $\theta(x)=\omega_1(x-c+L)$, $w(x)=\cos\theta(x)$, and
$J=[w(x_+),w(x_-)]\subset(-1,1)$, define
$P_{0,\boldsymbol g}(w)=\sum_{j=1}^mja_jT_j(w)$. The threshold is
\begin{equation}\label{eq:positive-threshold}
 \tau_0(\boldsymbol g)=-\min_{x\in I}f_{\boldsymbol g}(x).
\end{equation}
\cref{thm:finite} applies with scalar coefficient $\gamma=a_0$.

The threshold is attained at a physical endpoint or a location whose
cosine coordinate is an interior real root in $J$ of
$P_{0,\boldsymbol g}$. This stationary polynomial has degree at most
$m$ when nonzero. All root multiplicities and reduced degrees are
included. If the stationary polynomial vanishes identically, the
curve $f_{\boldsymbol g}$ is constant and either endpoint suffices.
\end{theoremEnd}

\begin{proofEnd}
\CombinedProofCheck{cor:positive}
Use the finite Gaussian model in \cref{sec:finite-model} and its compact
physical interval. Separating the intercept is an invertible permutation
of the coefficients, with $\min_Ih_0=1$, so \cref{thm:finite} at
order zero gives \cref{eq:positive-threshold} and the conditioned law.
To locate the minimum over the whole interval, use the Chebyshev
identities to write
\[
 f_{\boldsymbol g}(x)=\frac{\sqrt{1-w(x)^2}}{\sqrt L}
             \sum_{j=1}^ma_jU_{j-1}(w(x)),
 \qquad
 \frac{\mathrm{d}f_{\boldsymbol g}}{\mathrm{d}x}(x)=
             \frac{\omega_1}{\sqrt L}P_{0,\boldsymbol g}(w(x)).
\]
Because the cosine is strictly decreasing for $0<\theta(x)<\pi$,
interior stationary locations correspond exactly to interior roots of
the displayed polynomial. The extreme value theorem and Fermat's
theorem therefore give the endpoint and stationary candidates, with
finitely many distinct roots, including repeated roots, when the
polynomial is nonzero and has degree at most $m$. If it vanishes
identically, the derivative is zero throughout $I$, so the mean value
theorem gives constancy. Although the intercept is
independent of the spectral coefficients in the prior, reconstruction
uses its conditional posterior given $\boldsymbol g$ and the data.
\end{proofEnd}

\begin{theoremEnd}[combined main,category=monotone]{corollary}[Monotonicity in the Laplacian expansion]
\CombinedResultCheck{cor:monotone}{corollary}\label{cor:monotone}
Under the finite Gaussian model in \cref{sec:finite-model}, suppose
$m\geq4$ and $I\subset[c-L/2,c+L/2]$. Set
$h_1=-\phi_{2,L}-\phi_{4,L}/4$,
$\boldsymbol g=(g_0,g_1,g_3,\ldots,g_m)^T$, and
$f_{\boldsymbol g}=g_0+\sum_{j=1}^mg_j\phi_{j,L}$, where
$g_2=0$, $g_4=a_4-a_2/4$, and $g_j=a_j$ for $j\notin\{2,4\}$.
Then $f_{\boldsymbol a}=f_{\boldsymbol g}+(-a_2)h_1$, with inverse
obtained by reinserting $a_2$ and setting $a_4=g_4+a_2/4$.

For $\theta(x)=\omega_1(x-c+L)$, $w(x)=\cos\theta(x)$, and
$J=[w(x_+),w(x_-)]\subset[-1/\sqrt2,1/\sqrt2]$, define
\begin{equation}\label{eq:monotone-polynomials}
 P_{\boldsymbol g}(w)=\sum_{\substack{1\leq j\leq m\\j\ne2}}jg_jT_j(w),
 \qquad D(w)=1+4w^2-8w^4.
\end{equation}
Since $D\geq1$ on $J$, the threshold is
\begin{equation}\label{eq:monotone-threshold}
 \tau_1(\boldsymbol g)=\max_{w\in J}\frac{-P_{\boldsymbol g}(w)}{D(w)}.
\end{equation}
\cref{thm:finite} applies with scalar coefficient $\gamma=-a_2$.

The threshold is attained at a physical endpoint or a location whose
cosine coordinate is an interior real root in $J$ of
$\frac{\mathrm{d}P_{\boldsymbol g}}{\mathrm{d}w}D-P_{\boldsymbol g}\frac{\mathrm{d}D}{\mathrm{d}w}$.
This stationary polynomial has degree at most $m+3$ when nonzero.
All root multiplicities and reduced degrees are included. If the
stationary polynomial vanishes identically, the quotient
$-P_{\boldsymbol g}/D$ is constant and either endpoint suffices.
A nonzero constant stationary polynomial gives only endpoint candidates.
\end{theoremEnd}

\begin{proofEnd}
\CombinedProofCheck{cor:monotone}
Use the finite Gaussian model in \cref{sec:finite-model} with the rank and closed padding interval in
\cref{cor:monotone}. Substitution of the inverse map proves the curve
identity and invertibility, while the identity
$T_j(w(x))=\cos(j\theta(x))$ gives the derivative formulas
\[
 \frac{\mathrm{d}f_{\boldsymbol g}}{\mathrm{d}x}(x)=\frac{\omega_1}{\sqrt L}P_{\boldsymbol g}(w(x)),
 \qquad \frac{\mathrm{d}h_1}{\mathrm{d}x}(x)=\frac{\omega_1}{\sqrt L}D(w(x)).
\]
The assumed padding confines $\theta(x)$ to $[\pi/4,3\pi/4]$, where
the cosine is strictly decreasing and
$D(w)=1+4w^2(1-2w^2)\geq1$, including at the physical endpoints.
The positive derivative bound therefore allows \cref{thm:finite}
to supply the conditioned law and strictness.
Because the map transports the original Gaussian covariance,
$g_4$ and $a_2$ retain their induced dependence.

Since $D>0$ on compact $J$, the extreme value theorem and Fermat's
theorem locate the maximum of $-P_{\boldsymbol g}/D$ at an endpoint
or an interior stationary point. The quotient rule gives the
stationary polynomial in the statement, whose degree is at most
$m+3$ because $P_{\boldsymbol g}$ and $D$ have degrees at most
$m$ and $4$. If this polynomial vanishes identically, the quotient
is constant by the mean value theorem and either endpoint suffices.
Otherwise it has finitely many distinct roots, each a candidate
when it lies in the interior of $J$, regardless of multiplicity.
This also covers reduced degrees, including a nonzero constant
stationary polynomial with no interior candidates.
\end{proofEnd}

\begin{theoremEnd}[combined main,category=convex]{corollary}[Convexity in the Laplacian expansion]
\CombinedResultCheck{cor:convex}{corollary}\label{cor:convex}
Under the finite Gaussian model in \cref{sec:finite-model}, suppose
$m\geq2$ and $I\subset(c-L,c+L)$. Set $h_2=-\phi_{1,L}$,
$\boldsymbol g=(a_0,a_2,\ldots,a_m)^T$, and
$f_{\boldsymbol g}=a_0+\sum_{j=2}^ma_j\phi_{j,L}$.
Then $f_{\boldsymbol a}=f_{\boldsymbol g}+(-a_1)h_2$, with inverse
obtained by reinserting $a_1$.

For $\theta(x)=\omega_1(x-c+L)$, $w(x)=\cos\theta(x)$, and
$J=[w(x_+),w(x_-)]\subset(-1,1)$, define the polynomial and threshold
\begin{equation}\label{eq:convex-threshold}
 Q_{\boldsymbol g}(w)=\sum_{j=2}^mj^2a_jU_{j-1}(w),
 \qquad \tau_2(\boldsymbol g)=\max_{w\in J}Q_{\boldsymbol g}(w).
\end{equation}
\cref{thm:finite} applies with scalar coefficient $\gamma=-a_1$.

The threshold is attained at a physical endpoint or a location whose
cosine coordinate is an interior real root in $J$ of
$\frac{\mathrm{d}Q_{\boldsymbol g}}{\mathrm{d}w}$. The polynomial $Q_{\boldsymbol g}$
has degree at most $m-1$, so its stationary polynomial has degree at
most $m-2$ when nonzero. All root multiplicities and reduced degrees
are included. If the stationary polynomial vanishes identically,
$Q_{\boldsymbol g}$ is constant and either endpoint suffices.
A nonzero constant stationary polynomial gives only endpoint candidates.
\end{theoremEnd}

\begin{proofEnd}
\CombinedProofCheck{cor:convex}
Use the finite Gaussian model in \cref{sec:finite-model} with the rank and open padding interval in
\cref{cor:convex}. The coordinate map is an invertible signed
permutation. Since $0<\theta(x)<\pi$ on compact $I$,
\[
 \frac{\mathrm{d}^2h_2}{\mathrm{d}x^2}(x)=\frac{\omega_1^2}{\sqrt L}\sin\theta(x)>0,
 \qquad
 \frac{\mathrm{d}^2f_{\boldsymbol g}}{\mathrm{d}x^2}(x)=
 -\frac{\omega_1^2}{\sqrt L}\sin\theta(x)\,Q_{\boldsymbol g}(w(x)).
\]
The second formula uses
$\sin(j\theta(x))=\sin\theta(x)\,U_{j-1}(w(x))$, while compactness
keeps the positive curvature in the first formula bounded away from
zero. Thus \cref{thm:finite} gives the claimed threshold and posterior
under the stated open padding condition.

The extreme value theorem and Fermat's theorem locate the maximum
of $Q_{\boldsymbol g}$ at an endpoint or an interior zero of
$\frac{\mathrm{d}Q_{\boldsymbol g}}{\mathrm{d}w}$. The degree of $Q_{\boldsymbol g}$
is at most $m-1$, so this derivative has degree at most $m-2$ when
nonzero. If it vanishes identically, $Q_{\boldsymbol g}$ is constant
and either endpoint suffices. Otherwise its finitely many distinct
interior roots enter the candidate set regardless of multiplicity,
with none when the derivative is a nonzero constant. The same
argument covers every reduced degree.
\end{proofEnd}
